\documentclass[11pt]{article}

\usepackage{series}
\usepackage{amsmath,amssymb,amsthm}
\usepackage{hyperref}

\newcommand{\Tr}{\operatorname{Tr}}
\newcommand{\id}{\mathbf{1}}
\newcommand{\cI}{\mathcal{I}}
\newcommand{\cL}{\mathcal{L}}
\newcommand{\cM}{\mathcal{M}}
\newcommand{\cC}{\mathcal{C}}
\newcommand{\cE}{\mathcal{E}}
\newcommand{\cO}{\mathcal{O}}
\newcommand{\cH}{\mathcal{H}}

\title{Contextuality and Sequential Measurement Order:\\
\large Distinct Obstructions with a Shared MTT Interface}

\author{Peter Nero}

\date{July 2026, Version 2}

\begin{document}
\maketitle

\begin{abstract}
Quantum contextuality and sequential measurement order effects are related,
but they are not the same mathematical phenomenon. Kochen--Specker
contextuality is a static obstruction to gluing locally compatible value
assignments into one global noncontextual assignment. An order effect is a
dynamical statement about the noncommuting composition of quantum instruments,
including their state-update maps. We formulate both structures explicitly
and prove two separation results. First, projective qubit measurements exhibit
an exact order effect even though the original projective Kochen--Specker
obstruction does not apply in dimension two. Second, two instruments can have
the same effects and identical one-step outcome probabilities while producing
different sequential statistics. Effect or valuation data therefore cannot
determine order dependence. A common MTT carrier may nevertheless source both
structures. We state the required context-indexed source, instrument, measure,
and naturality data for such a unification and identify the missing comparison
theorem. The current q79 binary one-anchor recorder supplies an exact
stopped-output law for one restricted apparatus domain; it does not yet provide
the family of contexts needed for a contextuality theorem or a general
contextuality--order equivalence.
\end{abstract}

\section*{Version 2 Revision Note}
\addcontentsline{toc}{section}{Version 2 Revision Note}
\begin{description}
\item[Supersedes] \emph{Why Quantum Contextuality and Measurement Order
Dependence Are the Same Phenomenon: A Projection-Based Shadow Bridge in Modal
Triplet Theory}, version 1.
\item[Reason] The earlier central theorem identified failure of a common basin
refinement with sequential order dependence. Its proof did not establish
either implication. It also treated Kochen--Specker contextuality, generalized
operational contextuality, incompatible observables, and instrument
disturbance as interchangeable notions.
\item[Resolution] This version states the instrument algebra explicitly,
separates static valuation/gluing obstructions from dynamic composition, and
gives exact counterexamples to the claimed equivalence. It replaces the old
bridge theorem with a typed MTT construction contract.
\item[Retained result] Contextuality and order effects can be invariants of one
richer context-indexed physical model. Projection, admissibility, stable
records, and a shared upper carrier remain plausible ingredients of that
model.
\item[Remaining boundary] MTT has not yet constructed a selected family of
physical apparatus instruments together with a contextuality scenario and a
theorem relating its global-section obstruction to instrument composition.
The canonical q79 binary recorder is exact only on its declared one-context
domain.
\end{description}

\tableofcontents

\section{The corrected relation}

Four notions were compressed in the earlier paper:
\begin{enumerate}
\item incompatibility of observables;
\item Kochen--Specker failure of a noncontextual global valuation;
\item generalized contextuality of preparations, transformations, or effects;
\item dependence of sequential statistics on instrument order.
\end{enumerate}
They can interact, but they use different data and answer different questions.
The corrected relation is
\[
\boxed{
\begin{gathered}
\text{one context-indexed physical carrier}\\
\Downarrow\\
\text{static contextuality data}
\qquad\text{and}\qquad
\text{dynamic instrument data}.
\end{gathered}}
\]
The two lower structures may share a source without being equal. A theorem
identifying them would need a map between their mathematical categories and
proof that the relevant obstructions correspond.

\begin{description}
\item[Basic data.] Contextuality uses contexts, compatible outcomes, and
empirical distributions. Sequential order uses completely positive instrument
maps.
\item[Question.] Contextuality asks whether local assignments can be glued
globally. Sequential order asks whether physical composition depends on order.
\item[Typical obstruction.] The former is the absence of a global section or
noncontextual model. The latter is noncommuting state update.
\item[Role of effects.] Effects are central to many contextuality scenarios,
but they are insufficient to determine a sequential order effect.
\end{description}

\section{Measurement as an ordinary physical instrument}

Measurement is a physical interaction. It does not acquire a special logical
status because an observer reads the record. The appropriate operational
object includes both outcome probabilities and the state made available to
later interactions.

\begin{definition}[Finite quantum instrument]
Let \(\cH\) be a complex Hilbert space. A finite quantum instrument
\(\cI^A=\{\cI_i^A\}_{i\in I_A}\) is a family of completely positive,
trace-nonincreasing maps on trace-class operators such that
\[
\sum_{i\in I_A}\cI_i^A
\]
is trace preserving. Its effects are
\[
E_i^A=\cI_i^{A*}(\id),\qquad
E_i^A\geq0,\qquad
\sum_i E_i^A=\id .
\]
\end{definition}

For an input density operator \(\rho\), the probability and conditional
post-measurement state are
\begin{equation}
p(i|A,\rho)=\Tr\!\left[\cI_i^A(\rho)\right]
            =\Tr(\rho E_i^A),
\qquad
\rho_{i|A}=
\frac{\cI_i^A(\rho)}{p(i|A,\rho)}
\label{eq:instrument}
\end{equation}
when the denominator is nonzero. This is the standard instrument distinction:
the effects determine one-step probabilities, while the maps
\(\cI_i^A\) also determine what can happen next
\cite{DaviesLewis1970}.

For instruments \(A\) and \(B\), the sequential joint laws are
\begin{align}
p_{A\rightarrow B}(i,j|\rho)
 &=\Tr\!\left[\cI_j^B\!\left(\cI_i^A(\rho)\right)\right],
\label{eq:forward}\\
p_{B\rightarrow A}(j,i|\rho)
 &=\Tr\!\left[\cI_i^A\!\left(\cI_j^B(\rho)\right)\right].
\label{eq:reverse}
\end{align}

\begin{definition}[Sequential order effect]
The pair \(A,B\) has a sequential order effect on \(\rho\) if a declared
comparison of the operational laws in
\eqref{eq:forward}--\eqref{eq:reverse} differs. For example, one may compare
the probability of obtaining the labelled result \(+\) in both positions.
\end{definition}

An order effect is therefore not merely the statement that two operators fail
to commute. It is a statement about an input state, two fully specified
instruments, and a chosen comparison of their sequential records.

\section{Contextuality is a gluing obstruction}

A sharp contextuality scenario starts with a set \(X\) of measurements and a
cover \(\cC\) whose elements are jointly measurable contexts. If every
measurement has outcome set \(O\), a local assignment on
\(C\in\cC\) is an element of \(O^C\). An empirical model gives a compatible
probability distribution on the assignments in every context.

A deterministic noncontextual assignment is one global element of \(O^X\)
whose restriction to each context is allowed. More general noncontextual
models are convex mixtures of such assignments. Contextuality is the failure
of the relevant global model to exist. The sheaf formulation makes this
local-to-global obstruction explicit \cite{AbramskyBrandenburger2011}.

For projective quantum measurements, the Kochen--Specker theorem rules out a
global value assignment satisfying the functional relations among commuting
observables in Hilbert-space dimension at least three
\cite{KochenSpecker1967}. This is a theorem about compatible valuations. It
does not specify a state-update map for any physical apparatus.

Generalized operational contextuality is broader. It asks whether
operationally equivalent preparations, transformations, or measurement events
must have identical representations in an ontological model
\cite{Spekkens2005}. That framework can include transformation
contextuality and can apply to qubits. It remains important to name which
notion is being used; ``contextuality'' alone does not turn a static
Kochen--Specker scenario into a theorem about sequential instruments.

\section{Exact separation results}

\subsection{Order dependence without projective Kochen--Specker contextuality}

Let
\[
|0\rangle=
\begin{pmatrix}1\\0\end{pmatrix},
\qquad
|+\rangle=\frac{1}{\sqrt2}
\begin{pmatrix}1\\1\end{pmatrix},
\qquad
P_z=|0\rangle\langle0|,
\qquad
P_x=|+\rangle\langle+|.
\]
Take the initial state \(\rho=P_z\). Let \(Z\) and \(X\) be the binary
projective measurements in the \(z\)- and \(x\)-bases, with their Lueders
instruments
\[
\cI_i^Z(\rho)=P_i^Z\rho P_i^Z,
\qquad
\cI_j^X(\rho)=P_j^X\rho P_j^X .
\]

\begin{proposition}[Exact qubit order effect]
\label{prop:qubit-order}
For the events \(z+\) and \(x+\),
\[
p_{Z\rightarrow X}(z+,x+|\rho)=\frac12,
\qquad
p_{X\rightarrow Z}(x+,z+|\rho)=\frac14.
\]
Thus a two-dimensional projective system can exhibit a sequential order
effect even though the original projective Kochen--Specker theorem does not
apply in dimension two.
\end{proposition}

\begin{proof}
The \(z+\) outcome occurs first with probability one and leaves the state
\(P_z\). Hence
\[
p_{Z\rightarrow X}(z+,x+|\rho)
=\Tr(P_xP_zP_x)=|\langle+|0\rangle|^2=\frac12.
\]
In the reverse order, the \(x+\) outcome occurs with probability \(1/2\),
leaves the state \(P_x\), and the later \(z+\) outcome has conditional
probability \(1/2\). Therefore
\[
p_{X\rightarrow Z}(x+,z+|\rho)=\frac12\cdot\frac12=\frac14.
\]
\end{proof}

\begin{remark}
Proposition~\ref{prop:qubit-order} does not say that every generalized notion
of contextuality is absent for qubits. Spekkens contextuality can occur in
dimension two. The proposition addresses the specific equivalence claimed in
version 1: ordinary sequential order dependence does not imply the original
Kochen--Specker global-valuation obstruction.
\end{remark}

\subsection{The same effects can have different sequential laws}

The second separation is more basic. A POVM records only effects, whereas an
instrument also records state updates.

\begin{theorem}[Instrument underdetermination]
\label{thm:instrument-underdetermination}
There exist two instruments with identical effects and identical one-step
outcome probabilities for every input state, but with different statistics
under the same subsequent measurement.
\end{theorem}

\begin{proof}
On a qubit let \(P_0=|0\rangle\langle0|\) and
\(P_1=|1\rangle\langle1|\). Define the Lueders instrument
\[
\cL_i(\rho)=P_i\rho P_i
\]
and a measure-and-prepare instrument
\[
\cM_i(\rho)=\Tr(P_i\rho)\sigma_i,
\qquad
\sigma_0=|+\rangle\langle+|,
\quad
\sigma_1=|-\rangle\langle-|.
\]
Both families are completely positive, their sums are trace preserving, and
both have effects \(P_i\). Therefore
\[
\Tr[\cL_i(\rho)]=\Tr[\cM_i(\rho)]=\Tr(P_i\rho)
\]
for every \(\rho\).

Now take \(\rho=P_0\), retain the first outcome \(0\), and then measure
\(\{P_0,P_1\}\) again. The Lueders instrument gives
\[
\Tr[P_0\cL_0(P_0)]=1,
\]
whereas the measure-and-prepare instrument gives
\[
\Tr[P_0\cM_0(P_0)]
=\Tr(P_0|+\rangle\langle+|)=\frac12.
\]
Thus the same effects and one-step probabilities do not determine the
sequential law.
\end{proof}

\begin{corollary}[No bare contextuality--order equivalence]
\label{cor:no-bare-bridge}
No theorem formulated only in terms of effects, compatibility contexts, or
global valuations can determine general sequential order effects. Instrument
update maps or equivalent dynamical data are indispensable.
\end{corollary}

\begin{proof}
The data listed in the corollary cannot distinguish the two instruments in
Theorem~\ref{thm:instrument-underdetermination}, while a subsequent
measurement does distinguish them.
\end{proof}

This corollary identifies the exact flaw in the old basin-atlas bridge. Failure
of a common invariant refinement was used simultaneously as a valuation
obstruction and as a dynamical noncommutation statement. Those properties
belong to different structures unless a comparison map has been supplied and
proved faithful to both.

\section{A shared MTT interface}

The failure of equivalence does not make a common-source program empty. It
specifies what that program must construct.

\begin{definition}[Context-indexed MTT measurement package]
A context-indexed MTT measurement package consists of:
\begin{enumerate}
\item an upper physical state space \(\Omega\) with selected preparation laws
\(\mu_\rho\);
\item a compatibility cover \((X,\cC)\) and a presheaf of record assignments;
\item for every apparatus context \(C\), a physical coupling and record map
\(r_C:\Omega\to O^C\);
\item for every sequential apparatus \(A\), a quantum instrument
\(\cI^A=\{\cI_i^A\}\);
\item a context-indexed projection or intertwiner that derives each
\(\cI_i^A\) from the same upper dynamics;
\item pushforward identities
\[
(r_C)_*\mu_\rho=e_C
\]
for the contextual empirical distributions; and
\item composition identities showing that upper sequential evolution
descends to
\(\cI_j^B\circ\cI_i^A\).
\end{enumerate}
\end{definition}

The static and dynamic questions can then be asked on one package:
\begin{align*}
\mathfrak{o}_{\mathrm{ctx}}
&=\text{obstruction to a compatible global section},\\
\mathfrak{o}_{\mathrm{ord}}(A,B,\rho)
&=\text{difference between declared sequential laws}.
\end{align*}
They are two invariants of one object. They are not automatically the same
invariant.

\subsection{What a genuine bridge theorem would require}

A nontrivial contextuality--order bridge must add a rule
\[
\Phi:
\{\text{global-section obstructions}\}
\longrightarrow
\{\text{instrument-composition obstructions}\}
\]
or a functor between the underlying categories, and prove at least:
\begin{enumerate}
\item \emph{typing}: \(\Phi\) is defined for the selected physical context
family, not for an analogy between labels;
\item \emph{naturality}: restrictions of contexts commute with the MTT
projection and with instrument composition;
\item \emph{measure preservation}: empirical supports and probabilities are
the pushforwards of one selected source law;
\item \emph{soundness}: a static obstruction mapped by \(\Phi\) produces the
declared dynamic obstruction;
\item \emph{completeness}, if equivalence is claimed: every such dynamic
obstruction comes from the static one.
\end{enumerate}

Proposition~\ref{prop:qubit-order} shows that completeness fails for the broad
class of all sequential instruments. Theorem
\ref{thm:instrument-underdetermination} shows that soundness cannot be obtained
from effects alone. A future bridge must therefore use a restricted,
geometry-selected instrument family and state its domain.

\section{What projection and fixed points may contribute}

The useful MTT intuition can now be stated without overclaim.

\subsection{Projection can expose different lower structures}

One upper evolution may be read through several context-indexed projections.
Its one-step record supports may define a compatibility scenario, while its
conditioned evolution may define an instrument. If the projections are
derived from one carrier, both lower descriptions have a common physical
origin.

Common origin is already valuable. It can explain why the same apparatus
labels occur in contextuality tests and sequential protocols. It may also
constrain which instruments are physically realizable. It does not prove that
the global-section obstruction equals noncommutativity of those instruments.

\subsection{Fixed points can support records}

Stable fixed points or attracting regions can model persistent apparatus
records after a coupling. For a context \(C\), an MTT construction could assign
record regions \(B_i^C\subset\Omega\) and require
\[
\mu_\rho(B_i^C)=\Tr(\rho E_i^C).
\]
This is a meaningful source equation. The regions, measure, and equality must
all be derived; normalized labels or noninvertible projection alone do not
produce the probability law.

Nor does a family of record basins automatically instantiate a
Kochen--Specker scenario. One must prove which basins represent the same event
when it appears in overlapping compatible contexts. That overlap or gluing
map is exactly where contextuality lives.

\subsection{Order belongs to the coupling, not only the partition}

For sequential experiments the physical maps must retain the conditioned
post-record state. A partition of state space can say where a record lies but
not, by itself, how a subsequent apparatus acts. The required MTT object is
therefore a context-indexed transition kernel or CP instrument, not merely a
basin atlas.

\section{Current MTT status}

\subsection{The exact restricted result}

The current q79 program has an exact result on a declared canonical domain.
For its binary one-anchor nondemolition Fock recorder, the selected normal
state emits the stopped output measure, and second-moment capture descent is
exact. No separate Born axiom, fitted probability, observed probability, or
stochastic primitive is added on that domain.

This closes a real instrument-and-measure statement for one binary apparatus.
It improves the old paper's unsupported claim that projection alone supplies
Born basin measures.

\subsection{What the restricted result does not establish}

One binary context cannot display a Kochen--Specker gluing obstruction. The
current result does not yet construct:
\begin{itemize}
\item a family of overlapping compatible apparatus contexts;
\item context-independent identification of shared measurement events;
\item every corresponding physical instrument from one selected geometry;
\item arbitrary-context Born descent;
\item controlled finite-bandwidth or non-Markov corrections;
\item or an objective rule selecting one ontic history.
\end{itemize}

The current status is therefore
\[
\begin{array}{ll}
\text{canonical binary q79 record law:}&\text{exact on its domain},\\
\text{general selected apparatus family:}&\text{open},\\
\text{MTT contextuality realization:}&\text{open},\\
\text{contextuality--order equivalence:}&\text{not established}.
\end{array}
\]

\section{Interpretive consequences}

\subsection{Contextuality does not force anti-realism}

Kochen--Specker excludes a particular kind of global noncontextual value
assignment. It does not prove that physical records are unreal, that
observers create reality, or that every realist theory is impossible.
Generalized contextuality similarly constrains ontological representations
under explicit operational-equivalence assumptions. MTT may seek a realist
upper carrier, but it must reproduce those no-go results rather than dissolve
them by renaming context-dependent variables.

\subsection{Disturbance is not a dismissive explanation}

Calling an order effect ``disturbance'' is not an explanation until the
instrument is specified. Equations
\eqref{eq:forward}--\eqref{eq:reverse} make the mechanism testable. Two
apparatus implementations of the same POVM can disturb differently, as
Theorem~\ref{thm:instrument-underdetermination} shows.

\subsection{No privileged measurement event is required}

The formalism treats preparation, coupling, record formation, and later
coupling as ordinary dynamics. ``Measurement'' names a physical process with a
classical record interface. The distinction between contextuality and order
does not add an observer postulate; it simply keeps the static and dynamic
parts of that process typed correctly.

\section{Completion program}

The strongest next theorem is not another verbal shadow bridge. It is the
construction of one finite, selected test case:
\begin{enumerate}
\item choose a contextuality scenario with overlapping contexts and an
experimentally meaningful inequality or logical support obstruction;
\item construct each apparatus coupling and instrument from the same q79
source rather than importing arbitrary CP maps;
\item prove the source-measure pushforwards for all contexts;
\item compute both the contextuality witness and every relevant sequential
law from that package; and
\item test whether a restricted natural transformation relates the two
obstructions, recording counterexamples if it does not.
\end{enumerate}

A qutrit Kochen--Specker or KCBS-type context family would test the static
side. A matched set of sequential instruments would test the dynamic side.
The output should be a table distinguishing:
\[
\text{same source},\quad
\text{correlated invariants},\quad
\text{one-way implication},\quad
\text{equivalence}.
\]
Only the last entry would justify the original title.

\section{Conclusion}

Quantum contextuality and sequential measurement order are not one theorem in
two guises. Contextuality concerns the impossibility of a noncontextual global
model for locally compatible data. Order dependence concerns the composition
of physical instruments. Qubit Lueders measurements already refute the broad
equivalence, and instruments with identical effects show why effect or
valuation data cannot recover the missing dynamics.

The MTT common-source idea survives in a sharper form. A selected upper carrier
may generate both a contextual empirical model and a family of sequential
instruments. Fixed points may provide stable records, and projection may
relate upper dynamics to lower apparatus maps. The required source,
intertwining, pushforward, and gluing theorems are now explicit. The current
q79 binary recorder closes one important local part of that construction, but
the multi-context bridge remains open.

% BEGIN MTT MANAGED COMPUTATIONAL EVIDENCE
\section*{Computational Evidence and Reproducibility}
The contextuality and sequential-order obstructions are established from their respective operational data. The open strict-upgrade ledger proves neither obstruction and is included only to delimit stronger claims of complete physical reconstruction.

The referenced rows are frozen to the curated results repository state identified below. Hashes are grouped in eight-character blocks for line breaking.
\begin{quote}\small
\textbf{Repository:} \url{https://github.com/PeterNero/mtt-results-repro}\\
\textbf{Commit:} \texttt{31247ebb 5c22f3fb b5443024 365433c6 ee0bff4a}\\
\textbf{Manifest:} \path{release/result_manifest.json}\\
\textbf{Manifest SHA-256:}\\
\texttt{fb399689 60b00584 631dbf53 1a708e18}\\
\texttt{ef928d6b 6d935119 c185d7f6 32b1e7cd}
\end{quote}

Tier labels are quoted verbatim from that manifest. A row used directly supports only the specific computational statement identified above; a corpus-state cross-check does not prove this paper's local theorems; and an open row is evidence of an unresolved obligation, never of closure.

\par\begingroup\dimen0=\pagegoal\advance\dimen0 by -\pagetotal\ifdim\dimen0<7\baselineskip\newpage\fi\endgroup
\paragraph{Open boundary (not evidence of closure).}
\begin{itemize}
\setlength{\itemsep}{0.25em}
\item \path{A05/strict_upgrade_ledger} (\emph{open}).\par Current 2/9 strict no-knob upgrade ledger.
\end{itemize}

No imported row changes theorem ownership or promotes a neighboring claim: all local statements retain their stated hypotheses, domains, and limitations.
% END MTT MANAGED COMPUTATIONAL EVIDENCE

\begin{thebibliography}{99}

\bibitem{KochenSpecker1967}
S.~Kochen and E.~P. Specker,
\emph{The Problem of Hidden Variables in Quantum Mechanics},
Journal of Mathematics and Mechanics \textbf{17} (1967) 59--87,
doi:10.1512/iumj.1968.17.17004.

\bibitem{DaviesLewis1970}
E.~B. Davies and J.~T. Lewis,
\emph{An Operational Approach to Quantum Probability},
Communications in Mathematical Physics \textbf{17} (1970) 239--260,
doi:10.1007/BF01647093.

\bibitem{Spekkens2005}
R.~W. Spekkens,
\emph{Contextuality for Preparations, Transformations, and Unsharp
Measurements},
Physical Review A \textbf{71} (2005) 052108,
doi:10.1103/PhysRevA.71.052108,
arXiv:quant-ph/0406166.

\bibitem{AbramskyBrandenburger2011}
S.~Abramsky and A.~Brandenburger,
\emph{The Sheaf-Theoretic Structure of Non-Locality and Contextuality},
New Journal of Physics \textbf{13} (2011) 113036,
doi:10.1088/1367-2630/13/11/113036,
arXiv:1102.0264.

\bibitem{NeroBornClassical2026}
P.~Nero,
\emph{Born-Compatible Record Measures and the Classical Concentration Limit:
Separate Theorems and Their MTT Interface},
version 2, MTT research manuscript, July 2026.

\bibitem{MTTResults}
P.~Nero,
\emph{MTT Results Reproducibility Repository},
\url{https://github.com/PeterNero/mtt-results-repro}.

\end{thebibliography}

\end{document}
